\appendix
\FloatBarrier
\section{Glossaire de lecture}\label{sec:glossaire}
Les entrées suivent l'ordre de première apparition dans le corps du
texte, après le résumé. Les symboles introduits ensemble sont regroupés ;
leurs définitions opératoires restent celles des sections correspondantes.

\begin{small}
\begin{longtable}{@{}>{\raggedright\arraybackslash}p{.27\linewidth}p{.66\linewidth}@{}}
\toprule Terme ou notation & Sens dans cet article \\
\midrule\endhead
Effet rebond & Écart à l'économie de consommation attendue après un gain
d'efficacité. Ici, l'expérience mesure une réponse productive non
proportionnelle, sans calibration d'une consommation finale. \\
Endogène ; émergence & Produit par les interactions et l'évolution du
système, et non imposé comme cible macroscopique. \\
Auto-organisation critique (SOC) & Hypothèse d'un régime se maintenant
spontanément près d'un seuil critique. Elle n'est pas démontrée ici. \\
Entité & Agent générique producteur, susceptible de prêter, d'emprunter
et de disparaître ; aucun de ces rôles n'est prédéfini. \\
$K_i$ ; stock productif & Capital physique mobilisable de l'entité $i$,
en joules du modèle ; distinct de sa valeur nette. \\
Technologie $f_i(K)=A_iK^{\gamma_i}$ & Production ajoutée au stock
pendant un pas ; $A_i$ est le coefficient d'extraction, $\gamma_i$
l'exposant de production. \\
$K^\circ$, $\gamma^\circ$ & Dotation de naissance et exposant de
production à la naissance dans le régime de référence homogène. \\
Principal $q$ ; service $rq$ & Montant nominal d'un prêt et intérêts
exigibles par pas ; le principal n'est pas amorti. \\
$C_i$, $D_i$, $NW_i$ & Créances, dettes et valeur nette
(\emph{net worth}) : $NW_i=K_i+C_i-D_i$. \\
$\lambda$, $\delta$, $\sigma$, $\rho$ & Intensité de naissance,
fraction dépréciée par pas, amplitude du choc log-normal et intensité
des rencontres de marché. \\
Institution de crédit & Règles fixant le sens, le montant et le taux
des prêts ; elle ne désigne pas ici une banque particulière. \\
Invariant ; bilan réel & Relation imposée par les opérations du modèle.
Le bilan réel suit les apports, la production et les pertes du stock. \\
$B_t$, $\eta_i$ & Nombre de naissances au pas $t$ et choc logarithmique
individuel, renouvelé à chaque pas. \\
Concavité & Diminution du rendement marginal avec le stock, pour
$0<\gamma<1$ ; origine du gain possible de réallocation. \\
Défaut de service & Incapacité à payer intégralement les intérêts dus
au pas considéré. \\
Stock ; flux ; $\Delta t$ & Quantité détenue ; quantité transférée ou
produite pendant un pas ; durée physique éventuelle de ce pas. \\
Exergie & Capacité d'une ressource à fournir du travail relativement à
un environnement de référence ; sa correspondance avec les droits
économiques reste ici une hypothèse. \\
$K_{\rm aut}$, $x_i$ & Stock fixe déterministe d'une entité isolée
sans bruit ni crédit ; capital réduit $x_i=K_i/K_{\rm aut}$. \\
$\kappa^\circ$, $a$ & Dotation réduite $K^\circ/K_{\rm aut}$ ;
coefficient alternatif $a=A(K^\circ)^{\gamma-1}$. \\
Rééchelonnement ; compensation & Le premier transforme tout l'état et
ses unités ; la seconde adapte ici la dotation des futures entrantes
pour conserver $\kappa^\circ$ après la hausse de $A$. \\
Tension $T$ ; $K_{\rm eq}$ & Rapport $K_{\rm aut}/K_{\rm eq}$ ;
stock fictif donnant la production moyenne par productrice.
Dans les protocoles, $T$ peut aussi désigner l'horizon d'une simulation. \\
$Y$, $N_{\rm prod}$ & Production totale du pas et nombre d'entités
présentes à la phase productive. \\
Rendement marginal $f'(K)$ & Production supplémentaire à la marge
par unité de stock supplémentaire. \\
$L$, $\Delta$, $p$ & Perte productive de la prêteuse, surplus productif
joint et part de ce surplus attribuée à la prêteuse : $rq=L+p\Delta$. \\
Faillite ; cascade ; avalanche & Suppression d'une entité en défaut
ou insolvable ; propagation des pertes ; composante faiblement connexe
des pertes entre mortes du même pas. \\
Racine ; génération ; $b_1$, $b_2$ & Mort initiale ; rang de propagation
intra-pas ; fraction des mortes non racines et nombre moyen de liens
causals entre générations successives par morte. \\
Graine ; appariement ; IC95 & Initialisation du générateur aléatoire ;
branches issues d'un même état ; intervalle de confiance à 95\,\%
calculé entre répétitions, sauf indication contraire. \\
$k$ ; stationnarité & Taille historique du bassin de rencontre, fixée
à deux dans le protocole principal ; invariance temporelle des propriétés
statistiques, non démontrée par un seul test de dérive. \\
Loi de Little & Relation entre effectif moyen, flux d'entrée et durée
de vie moyenne sous conditions de stabilité et de maîtrise des bords. \\
$S$, $\alpha$, $S_c$, $Z$ & Taille d'avalanche, exposant, coupure
exponentielle et constante de normalisation de sa loi ajustée. \\
Bootstrap ; $R^2$ & Rééchantillonnage de réalisations entières ;
score descriptif d'ajustement, non preuve d'une famille de lois. \\
Contraction ; $D$, $\beta$ & Suite de pas de baisse d'une grandeur ;
durée de cette suite et taux de décroissance de sa survie exponentielle.
$D$ sans indice dans ce contexte n'est pas une dette. \\
Gini ; Lorenz & Indice synthétique d'inégalité ; courbe de part cumulée
d'une grandeur en fonction de la part d'entités classées par elle. \\
Queue ; Pareto ; log-normale & Domaine des grandes valeurs ; famille
à décroissance en puissance ; famille dont le logarithme est normal. \\
Élasticité finie $\varepsilon$ & Rapport du logarithme de la réponse
productive au logarithme du facteur technique $m$, ici $m=1{,}5$. \\
Quantile ; décile ; cohorte & Seuil de classement ; groupe de 10\,\%
reclassé à chaque date ; ensemble d'individus suivis par identité. \\
Covariance & Mesure de variation conjointe utilisée dans l'identité
de pondération du partage ; elle ne démontre pas une causalité. \\
Taux de charge & Intérêts versés ou dus, selon la définition,
rapportés à la production plus les intérêts reçus ; distinct d'un
rapport stock de dette/production. \\
Vuong & Comparaison de vraisemblances de modèles sur un même
échantillon ; ses hypothèses et ses limites sont explicitées en annexe. \\
\bottomrule
\end{longtable}
\end{small}

\clearpage
\section{Vérifier le régime étudié et les calculs de réponse}\label{sec:stationnarite}
Pour chaque graine $s$, bras et observable $X$, nous calculons
\[
 R_s(X)=\frac{\langle X\rangle_{3750<t\le4000}}
                  {\langle X\rangle_{3000<t\le3250}},\qquad
 t_X=\frac{\overline R(X)-1}{s_R/\sqrt{12}},
\]
où $s_R$ est l'écart-type des douze rapports. L'IC95 est
$\overline R\pm t_{0{,}975;11}s_R/\sqrt{12}$, avec
$t_{0{,}975;11}\simeq2{,}201$.
Ce diagnostic compare deux niveaux temporels éloignés ; il n'analyse
ni toute la structure de corrélation ni la stabilité de la distribution.

\begin{table}[htbp]\centering\small
\caption{Rapport des moyennes des derniers et des premiers 250 pas de la fenêtre terminale.}\label{tab:stationnarite}
\begin{tabular}{@{}lrrr@{}}
\toprule
Bras & $R(Y)$ (IC95) & $R(N)$ (IC95) & $R(\sum K)$ (IC95) \\
\midrule
Contrôle & $1{,}0029\pm0{,}0152$ & $1{,}0014\pm0{,}0132$ & $1{,}0043\pm0{,}0178$ \\
$A\times1{,}5$ & $0{,}9968\pm0{,}0139$ & $0{,}9979\pm0{,}0116$ & $0{,}9957\pm0{,}0167$ \\
$A\times1{,}5$, dotation compensée & $0{,}9987\pm0{,}0132$ & $0{,}9974\pm0{,}0122$ & $0{,}9998\pm0{,}0148$ \\
Entrantes $A\times1{,}5$ & $1{,}0088\pm0{,}0106$ & $1{,}0070\pm0{,}0090$ & $1{,}0092\pm0{,}0128$ \\
Entrantes $A\times0{,}75$ & $0{,}9906\pm0{,}0159$ & $0{,}9928\pm0{,}0139$ & $0{,}9905\pm0{,}0184$ \\
Entrantes $\gamma=0{,}6$ & $1{,}0097\pm0{,}0159$ & $1{,}0038\pm0{,}0133$ & $1{,}0100\pm0{,}0186$ \\
\bottomrule\end{tabular}\end{table}


Aucun de ces contrastes n'excède le seuil bilatéral usuel, mais les
intervalles autorisent des dérives de l'ordre du pour-cent. Les observations
temporelles d'un run ne sont pas utilisées comme 1000 répétitions indépendantes.
Le résultat ne vaut pas automatiquement pour une fenêtre antérieure,
une autre observable ou une configuration exploratoire non présentée.

\paragraph{Renouvellement et mémoire de l'état initial.}
Le renouvellement d'un décile fournit une autre échelle temporelle, mais
il faut préciser la grandeur de classement et distinguer la disparition
d'une entité de sa sortie du décile. Pour un ensemble $D(t_0)$ défini
à une date de référence, la rétention
$|D(t_0)\cap D(t_0+h)|/|D(t_0)|$ mesure la persistance dans le groupe,
tandis que la fraction de ses membres encore vivants mesure leur survie.
Ces deux courbes ne sont pas interchangeables ; le terme « premier
décile » doit préciser s'il désigne les valeurs basses ou les plus élevées.
Même une rétention faible ne démontre pas l'oubli de la configuration
initiale : des engagements et des corrélations peuvent persister au-delà
des individus. Ce contrôle doit être confronté aux autocorrélations des
observables et à des initialisations distinctes. Aucun temps d'oubli
quantifié n'est établi ici ; l'éloignement de la fenêtre finale et
l'absence de dérive détectée ne suffisent pas à le certifier.

Les bilans de population donnent exactement
$N(t+1)-N(t)=B(t)-D(t)$ si $B$ et $D$ désignent les entrées et sorties
du pas avec la même convention de comptage. Sur une longue fenêtre stable,
$\overline B\simeq\overline D$ et $\overline B\simeq\lambda$ ;
le rapport $\overline D/\overline N$ devient proche de
$\lambda/\overline N$. Ce n'est pas exactement la moyenne de
$D(t)/N(t)$, et les termes de bord ne sont pas nuls par décret.
La relation explique pourquoi une élasticité de mortalité peut être
presque l'opposée de celle de population sans constituer une découverte
causale indépendante.

L'élasticité de l'équation~\eqref{eq:eps} porte sur une intervention
finie. Sa moyenne et son IC sont calculés après formation des douze
log-rapports appariés, et non après moyenne de tous les points temporels
dans les deux bras. Les formules de décomposition sont sensibles au
choix de l'effectif : les entités présentes à la production peuvent
disparaître avant l'instantané de fin de pas. Cette différence est
conservée dans les données sources ; elle ne doit pas être effacée pour
forcer la fermeture numérique d'une identité mal définie.

\section{Ajustements de distributions : définitions et limites}\label{sec:ajustements}
\subsection{Masse, densité par classe et fonction de survie}
Pour une taille entière $S$, la masse $P(S=s)$ et la survie
$P(S\ge s)$ sont distinctes. L'histogramme représente ici la masse
d'une classe $[l,r[$ divisée par $r-l$, afin qu'une classe large ne
paraisse pas plus probable seulement parce qu'elle contient plus
de tailles. La courbe ajustée est sommée sur les mêmes tailles
entières avant cette division. Une pente de survie ne peut pas être
substituée directement à l'exposant de masse.

La vraisemblance descriptive de la loi~\eqref{eq:avalanche} est maximisée
par graine, sur toutes les tailles $s\ge1$. La normalisation numérique
utilise $1\le s\le20000$ ; la masse des mille dernières valeurs est
contrôlée inférieure à $10^{-10}$ à l'optimum. Les bornes numériques
de recherche sont $-2\le\alpha\le5$ et
$10^{-5}\le1/S_c\le2$ ; un optimum sur une borne signalerait une
identification insuffisante. Les 24 ajustements présentés sont intérieurs.
Le support d'ajustement n'est pas choisi pour maximiser la rectitude
apparente d'un sous-ensemble de points.

Les IC inter-graines et le bootstrap de runs entiers mesurent la
variabilité des estimateurs ou des courbes moyennes obtenus par ce
protocole. L'ajustement par produit de probabilités ne démontre pas
l'indépendance des avalanches d'un run. Nous n'en tirons donc pas
d'erreurs-types événementielles ni un test d'adéquation certifié.
Il faudrait notamment compléter par une comparaison de familles sur
le même support et un diagnostic d'adéquation absolue tenant compte
de la dépendance temporelle \cite{clauset}.

Le coefficient descriptif affiché est
\[
 R^2_{\log}=1-
 \frac{\sum_j[\log\widehat p_j-\log p_{j,\rm ajust}]^2}
      {\sum_j[\log\widehat p_j-\overline{\log\widehat p}]^2},
\]
sur les classes empiriques non vides. Ce n'est pas un pourcentage
de probabilité correctement prédit, et ce n'est pas le critère
optimisé par la vraisemblance. Il dépend du découpage des classes.
Les zéros ne sont pas remplacés par une petite valeur arbitraire
pour les rendre logarithmables.

\subsection{Durées : exponentielle discrète et choix de fenêtre}
Si $D\ge1$ suit une loi géométrique, avec probabilité d'arrêt $u$ à
chaque pas, alors $P(D\ge d)=(1-u)^{d-1}$ et
$\widehat u=1/\overline D$. Ainsi
$\widehat\beta=-\log(1-1/\overline D)$ dans
$P(D\ge d)=\exp[-\beta(d-1)]$.
Les régressions de la figure~\ref{fig:art_cycles} utilisent cette
forme normalisée, non une droite en log-log. Le $R^2$ y est calculé
sur les log-survies non nulles, aux durées 1 à 11. Les épisodes restent
potentiellement corrélés et les frontières de fenêtre tronquent certains
d'entre eux ; les IC sont donc fondés sur les cinq graines et non sur
un nombre fictif de récessions indépendantes.

\subsection{Le test de Vuong et ce qu'il compare réellement}
Dans les analyses existantes des valeurs nettes et intérêts, le seuil
$x_{\min}$ est choisi par minimisation d'une distance de
Kolmogorov--Smirnov pour une puissance continue. Au-dessus de ce même
seuil et sur les mêmes valeurs, on ajuste aussi une log-normale et
une exponentielle, conditionnées à $X\ge x_{\min}$.
La puissance comparée ici est \emph{sans coupure supérieure} : ce
n'est pas le modèle discret à coupure exponentielle des avalanches.
La troncature \emph{à gauche}, qui définit un échantillon de queue,
ne doit pas être confondue avec la coupure \emph{à droite}.

Pour les densités ajustées $f$ et $g$, on définit
$d_i=\log f(x_i)-\log g(x_i)$ puis
\[
 z=\frac{\sqrt n\,\overline d}{s_d},\qquad
 s_d^2=\frac1{n-1}\sum_i(d_i-\overline d)^2.
\]
Une valeur positive favorise $f$ en log-vraisemblance moyenne.
L'interprétation usuelle de $\pm1{,}96$ comme seuil à 5\,\%
nécessite des hypothèses de régularité, de séparation des modèles
et d'échantillonnage. Ici le seuil a été sélectionné, les entités
interagissent, et certaines log-normales ajustées sont très proches
d'une limite en puissance. Le score est donc rapporté comme diagnostic,
pas comme certification de ces hypothèses \cite{vuong}.

Le signalement $\sigma_{\rm LN}>10$ concerne 50 des 72 ajustements
d'intérêts et 46 des 72 ajustements de valeur nette. Il marque une
zone numérique, pas une frontière mathématique entre deux familles.
Deux groupes visibles dans un nuage de paramètres d'ajustement ne
signifient pas que les agents forment deux populations économiques :
chaque point de ce nuage représentait un ajustement, pas une entité.
Les résultats par ajustement et leurs identifiants sont disponibles dans
\chemin{data_revision/rev_vuong_points.csv}.

Comparer Pareto, puissance à coupure et log-normale demande de
fixer le même domaine, de décrire la sélection du seuil et d'évaluer
l'adéquation, pas seulement de compter les paramètres. Une coupure
qui divergerait avec la taille du système définirait une limite
intéressante, à vérifier sur plusieurs tailles ; elle ne rend pas
gratuit son ajustement à taille finie.

\section{Charge du crédit et pondération du surplus}
\subsection{Définition et mesure du taux de charge}\label{sec:charge}
Notons $P_i$ la production du pas, $I_i^{\text{reçu}}$ les intérêts
reçus, $I_i^{\text{versé}}$ ceux effectivement payés, et $D_i$ la
somme des principaux dus. Trois rapports doivent être distingués :
\[
 c_i^{\text{versé}}=\frac{I_i^{\text{versé}}}{P_i+I_i^{\text{reçu}}},\qquad
 c_i^{\text{dû}}=\frac{\sum_{\ell:i\,\mathrm{emprunte}}r_\ell q_\ell}
                         {P_i+I_i^{\text{reçu}}},\qquad
 d_i^{\text{revenu}}=\frac{D_i}{P_i+I_i^{\text{reçu}}}.
\]
Le premier est une part de revenu brut versée ; le deuxième est la
charge contractuelle due ; le troisième rapporte un stock de dette
à un revenu par pas et s'interprète en nombre de pas. La
variable $D_i/P_i$ ne représente aucun des deux premiers rapports.
Le dénominateur inclut les revenus financiers, même si leur
somme agrégée n'est pas une nouvelle production réelle.

\figurine{art_charge}{0.95}{\textbf{Un diagnostic de charge correctement
nommé.} Règle marginale M4.4, douze graines. Pour les instantanés
$3000<t\le3990$, la variable est
$I_i^{\text{versé}}/(P_i+I_i^{\text{reçu}})$ ; la réponse est la fraction
des entités décédant dans les dix pas suivants. Les classes ont les
bornes explicites indiquées, communes aux graines ; elles ne sont
pas des classes de capital. La dernière fenêtre incomplète est
exclue. Les IC95 comparent les fractions obtenues dans douze runs,
sans attribuer une indépendance binomiale aux entités répétées.
Les décès toutes causes et les racines par insolvabilité sont séparés.}

Les panneaux enregistrent les intérêts versés, pas le service dû
au moment exact précédant le paiement. Un paiement au prorata peut
faire paraître faible une charge pourtant insoutenable ; le moment
de prélèvement et les flux reçus dans le pas comptent aussi.
Le diagramme reste descriptif, conditionnel aux survivantes observées
dans le panneau. Il ne justifie ni une causalité de la charge vers
le décès ni une réfutation générale de convexité. Des calculs de
convexité ou d'écart de Jensen sur $D_i/P_i$ ne répondraient pas
à la question posée sur le taux de charge ainsi défini.

\subsection{Preuve de l'identité de pondération}\label{sec:ponderation}
Pour $n$ contrats dont le surplus moyen est strictement positif,
posons $\overline p=n^{-1}\sum p_j$,
$\overline\Delta=n^{-1}\sum\Delta_j$ et
$p_w=\sum p_j\Delta_j/\sum\Delta_j$.
En définissant la covariance empirique avec le dénominateur $n$,
\[
 \operatorname{Cov}_n(p,\Delta)
 =\frac1n\sum_j(p_j-\overline p)(\Delta_j-\overline\Delta)
 =\frac1n\sum_jp_j\Delta_j-\overline p\,\overline\Delta,
\]
on obtient immédiatement
\begin{equation}
 p_w-\overline p=\frac{\operatorname{Cov}_n(p,\Delta)}{\overline\Delta}.
 \label{eq:identite}
\end{equation}
Si une covariance corrigée avec dénominateur $n-1$ est utilisée,
il faut multiplier le numérateur par $(n-1)/n$. L'identité vaut
aussi en espérance lorsque les moments existent et que
$\mathbb E[\Delta]\ne0$.

Pour un partage fixé constant, l'écart est nul par construction.
Pour une règle variable, il peut être positif même si la plupart
des contrats ont un principal ou un surplus très petit. La
distribution effective des contrats décide de son ampleur ; une
grille analytique couvrant des contrats rarement réalisés ne peut
pas remplacer cette distribution. Qualifier une règle à partir
de ces seuls cas théoriques surestimerait leur portée empirique.
L'identité comptable ne démontre pas une explication des différences
de mortalité ou de population entre institutions.

\section{Compléments de mesure et valeurs tabulées}\label{sec:tableaux}
\subsection{Rapports de quantiles}
Le tableau ci-dessous et la figure~\ref{fig:art_quantiles} reposent
sur les mêmes observations. Chaque $\pm$ est un demi-IC95 de
Student des douze rapports par graine, non un écart-type entre
entités. « Revenu brut » signifie production plus intérêts reçus.
\begin{small}% Généré : rapports par graine, IC95 Student (11 ddl).
\begin{longtable}{@{}lrrrr@{}}
\toprule
Grandeur & q10 & q50 & q99 & q99,9 \\
\midrule
\endhead
Production & $1{,}748 \pm 0{,}004$ & $1{,}806 \pm 0{,}003$ & $1{,}831 \pm 0{,}004$ & $1{,}816 \pm 0{,}008$ \\
Revenu brut & $1{,}755 \pm 0{,}004$ & $1{,}776 \pm 0{,}007$ & $1{,}741 \pm 0{,}070$ & $1{,}474 \pm 0{,}060$ \\
Valeur nette & $1{,}103 \pm 0{,}007$ & $1{,}486 \pm 0{,}016$ & $1{,}452 \pm 0{,}095$ & $1{,}169 \pm 0{,}065$ \\
Capital & $1{,}386 \pm 0{,}008$ & $1{,}451 \pm 0{,}005$ & $1{,}489 \pm 0{,}006$ & $1{,}460 \pm 0{,}013$ \\
Intérêts reçus & $0{,}111 \pm 0{,}025$ & $1{,}743 \pm 0{,}013$ & $1{,}726 \pm 0{,}081$ & $1{,}447 \pm 0{,}062$ \\
\bottomrule
\end{longtable}
\end{small}

\subsection{Portée temporaire d'une intervention}
Une intervention limitée aux entités présentes peut s'éteindre
avec leur cohorte si les naissances ne l'héritent pas. Cette
expérience de portée partielle est un contrôle de protocole,
non une preuve principale du mécanisme de rebond. À l'inverse,
modifier toutes les entités et les entrantes futures conserve
le traitement malgré le renouvellement de population. Modifier
$\gamma$ change la forme de la fonction et les dimensions de
son coefficient : ce n'est pas la même opération que multiplier
$A$ à $\gamma$ fixe.

\subsection{Plusieurs pertes et causes suffisantes}
Le graphe peut attribuer plusieurs pertes à une victime.
Cela ne veut pas dire que chacune aurait suffi, isolément, à
la rendre insolvable. Pour l'illustrer, une valeur nette initiale
de 5 et deux pertes de 3 conduisent à $-1$ ; aucune perte seule
ne suffit. Avec une valeur nette de 2, chacune des deux pertes
suffirait. C'est la marge de solvabilité avant les pertes qui
fait la différence, pas le seul nombre d'arêtes.

Le comptage des parents doit aussi préciser s'il retient les
seules pertes venant de la génération précédente ou toutes
celles du pas. L'analyse détaillée de suffisance exige de
reconstituer également les dettes effacées ; elle n'est pas
nécessaire pour établir l'existence de la propagation ni la
non-proportionnalité de la réponse. Les calculs antérieurs
restent archivés, mais leurs pourcentages ne sont pas utilisés
comme preuve dans cet article.

\subsection{Stock et production sur une même réalisation}
\figurine{rev_stock_trajectoire}{0.95}{\textbf{Observer la différence entre
stock et flux.} M4B central, graine 11, extrait $3000<t\le3200$.
Chaque ligne suit une seule réalisation, sans incertitude de mesure
artificiellement ajoutée ; les points orange repèrent les pas
de baisse. Les durées de contraction sont calculées séparément
sur les deux observables dans la figure~\ref{fig:art_cycles}.}

\FloatBarrier
\section{Sources locales, configurations et reproductibilité}\label{sec:sources_revision}
Le PDF est compilable avec les seuls fichiers LaTeX et figures
de ce dossier, sans le rapport de stage ni les questionnaires.
Les calculs nouveaux sont produits par
\chemin{../scripts/revision_article.py}, qui n'écrit que dans
le dossier de cet article et ne lance aucune simulation.
Les exports numériques, observations par graine et empreintes
SHA-256 des sources sont dans \chemin{data_article/}.
Les compléments conservés de la version précédente sont dans
\chemin{data_revision/}. Recalculer depuis les microdonnées
demande les campagnes locales, qui ne sont pas toutes versionnées.

\subsection{Trois protocoles, sans fusion des échantillons}
\begin{small}
\begin{longtable}{@{}p{.18\linewidth}p{.75\linewidth}@{}}
\toprule Famille & Configuration et données \\\midrule\endhead
M4B central & $A=1$, $\gamma=1/2$, $\lambda=30$, $\delta=0{,}05$,
$\sigma=0{,}25$, $K^\circ=25$, $k=3$, cible géométrique. Cinq graines
11--15 ; contractions sur $1000<t\le4000$, Gini sur l'état final.
Les identifiants Simulation Lab sont dans les exports d'épisodes et de
confirmation. Le balayage du Gini présenté varie $\sigma$ seulement. \\
Recalage Live-v1 & Référence $A=1$, $\gamma=1/2$, $\lambda=30$,
$\delta=\sigma=0{,}01$, $K^\circ=25$, $\rho=1$ ; graines 0--2,
départ vide. Fenêtre $1200<t_{\rm ancien}\le2000$.
Le bras à dépréciation fine utilise sa référence $\delta=0{,}002$.
Les dossiers sont nommés dans \texttt{temps\_graines.csv}. \\
M4.4 principal & Même référence physique que Live-v1 ; cible arithmétique
en régime homogène, sens libre du prêt, règle marginale. Douze graines
0--11, intervention au pas 2001, fenêtre $3000<t\le4000$.
Dotation compensée : $K^\circ=56{,}25$ pour les nouvelles entrantes
après intervention ; stocks existants inchangés. \\
Partage M4.4 & $p\in\{0;0{,}25;0{,}5;0{,}75;1\}$ ou règle marginale ;
douze graines, même fenêtre. Les rapports au bras $p=0$ sont appariés.
Diagnostic de charge : règle marginale seulement, $3000<t\le3990$,
horizon de décès de dix pas. \\
Rencontres M4.4 & $\rho\in\{0{,}5;1;1{,}5;2;3\}$ ; douze graines
par valeur. Les Gini de bilan sont recalculés sur les panneaux de
$3000<t\le4000$, comme les fenêtres des autres mesures. \\
\bottomrule
\end{longtable}
\end{small}

\subsection{Chaque figure et ses fichiers de points}
Les fichiers CSV sans préfixe et \chemin{diagnostic_chaine.json} sont
relatifs à \chemin{data_article/}. Les fichiers TeX et les chemins
commençant par \chemin{data_revision/} sont relatifs au dossier LaTeX.
Le manifeste relie les exports aux sources
originales et à leurs empreintes ; il ne transforme pas un export
agrégé en microdonnées complètes.
\begin{footnotesize}
\begin{longtable}{@{}p{.09\linewidth}p{.48\linewidth}p{.36\linewidth}@{}}
\toprule Fig. & Données / fichier local & Nature du calcul \\
\midrule\endhead
\ref{fig:exemple} & \chemin{exemple_article.tex} & Quatre entités fictives,
prêt réel et inscription nominale ; aucune mesure simulée. \\
\ref{fig:art_technologie} & \chemin{technologie_points.csv},
\chemin{technologie_courbes_v2.csv} & Technologies analytiques
$2K^{1/4}$ et $K^{1/2}$ ; pas de régression. \\
\ref{fig:cascade} & \chemin{cascade_article.tex} & Bilan fictif à cinq
entités ; arbre des quatre faillites vérifié séparément. \\
\ref{fig:art_avalanches} & \chemin{avalanches_counts.csv},
\chemin{avalanches_fits.csv}, \chemin{avalanches_points.csv} & Tailles
entières dans les CSV d'avalanches ; ajustements et IC inter-graines. \\
\ref{fig:art_cycles} & \chemin{cycles_fits.csv}, \chemin{cycles_points.csv} ;
épisodes dans \chemin{data_revision/rev_cycles_episodes.csv} & Cinq
runs M4B ; coefficient exponentiel depuis la durée moyenne. \\
\ref{fig:art_lorenz} & \chemin{lorenz_points.csv}, \chemin{lorenz_ginis.csv} &
Panneaux du contrôle ; dates exactes 2010, 3000, 4000. \\
\ref{fig:art_sensibilite} & \chemin{m4b_sensibilite_points.csv},
\chemin{m4b_nw_confirmation.csv} & Gini NW sur cinq instantanés finaux
indépendants par cellule. \\
\ref{fig:art_rho} & \chemin{rho_points.csv}, \chemin{rho_fits.csv},
\chemin{rho_nw_graines.csv} & Trois estimateurs du tableau de campagne,
plus Gini NW recalculé ; régressions par graine. \\
\ref{fig:art_temps} & \chemin{temps_graines.csv}, \chemin{temps_points.csv} &
Séries Live-v1 ramenées à une même durée et à un même flux par pas ancien. \\
\ref{fig:art_reponse} & \chemin{elasticites_graines.csv} ; courbes dans
\chemin{data_revision/rev_reponse_courbes.csv} & Élasticités recalculées
depuis les séries ; lissage centré 101 pas conservé pour les courbes. \\
\ref{fig:chaine} & \chemin{corps_article.tex} ; diagnostics
\chemin{diagnostic_chaine.json} & Schéma interprétatif,
pas graphe causal identifié par intervention sur chaque médiateur. \\
\ref{fig:art_quantiles} & \chemin{data_revision/rev_quantiles_graines.csv} &
Rapports par graine conservés ; abscisse transformée en logit. \\
\ref{fig:art_deciles} & \chemin{deciles_v2_points.csv},
\chemin{deciles_v2_graines.csv} & Parts de production, intérêts et NW ;
classement par capital à chaque instantané. \\
\ref{fig:art_institutions} & \chemin{institutions_points.csv},
\chemin{institutions_fits.csv}, \chemin{institutions_references.csv} &
Rapports à $p=0$ et régressions linéaires descriptives. \\
\ref{fig:art_charge} & \chemin{charge_graines.csv}, \chemin{charge_points.csv},
\chemin{charge_resume.csv} & Intérêts versés/revenu brut ; décès
observables jusqu'au bout de l'horizon. \\
\ref{fig:rev_stock_trajectoire} &
\chemin{data_revision/rev_stock_trajectoire_points.csv} &
Extrait brut M4B ; figure conservée. \\
\bottomrule
\end{longtable}
\end{footnotesize}

Les nouveaux contrôles de composition par âge sont dans
\chemin{ages_graines.csv}. Les conventions de classement sont
explicites dans le script ; les ex æquo sont départagés par ordre
stable. Le plan de révision et le registre des décisions sont hors
du dossier LaTeX, dans \chemin{../revision_2026-09-15/}, pour
préserver séparément les corrections à reporter au rapport de stage.

Le complément sur les médianes instantanées des âges est calculé par
\chemin{../scripts/ages_v2.py}. Ses valeurs par graine, sa convention
de moyenne et les empreintes des panneaux sources sont conservées dans
\chemin{../revision_2026-09-17/ages_v2_graines.csv} et
\chemin{../revision_2026-09-17/ages_v2.json}. Il n'exécute aucune
simulation et ne modifie aucun générateur de figures.


Les configurations détaillées et les identifiants M4.4 de la
campagne initiale restent dans
\chemin{data_revision/runs_m44.csv}, qui distingue les paramètres
enregistrés et les interventions. Certaines métadonnées d'amorçage
ne renseignent pas les paramètres ; ces lacunes ne sont pas comblées
par déduction silencieuse.
Les codes techniques \texttt{K0} des fichiers sources sont
conservés pour leur traçabilité ; dans la notation scientifique,
la dotation se lit désormais $K^\circ$ et son rapport réduit
$\kappa^\circ$. Renommer la notation ne modifie pas les moteurs.

\subsection{Compatibilité des moteurs et changements de règles}
\label{sec:compatibilite}
Trois niveaux de contrôle doivent être distingués : conserver les
fonctionnalités du laboratoire de simulation, reproduire une trajectoire
dans un régime de référence, et comparer deux modèles dont les règles
diffèrent. Seul le deuxième porte directement sur la parité dynamique.

Le rapport M4.2, section « Parité M4B », documente des tests à
$A=1$, $\gamma=1/2$ et $k=2$ : identité des variables discrètes
et comparaison des flottants à tolérance relative $10^{-9}$.
La sonde de 1000 pas rapportée ne présente pas de divergence discrète,
avec un écart relatif maximal des séries de $3{,}1\,10^{-12}$.
Ce constat n'est pas une identité bit à bit ni une garantie à tout horizon.

Pour M4.4, dans le dossier de campagne
\chemin{m4_4_rebond_credit_soc/}, le fichier
\chemin{results/analysis/parity_deviations_8000.csv}
conserve 8000 écarts maximaux nuls sur les 26 colonnes comparées à
la réalisation M4.3 de référence, dans le régime homogène à cible
arithmétique. Le test \chemin{tests/test_parity_m4_3.py} définit ce
périmètre. Le test \chemin{tests/test_v1_equivalence.py} compare
séparément Live-v1 au mode de compatibilité
\texttt{richest\_lends}, avec une intervention hétérogène ; un contrôle
négatif vérifie que libérer le sens du prêt peut changer la trajectoire.
La présence de ce test définit une exigence vérifiable ; elle ne remplace
pas à elle seule un compte rendu d'exécution pour chaque version.

Les chemins de campagnes M4.4 sont conservés comme identifiants de
provenance. Le passage à une règle de marché différente doit être
distingué d'une extension d'instrumentation, sans renommer les archives
dont proviennent les données. La non-régression ne
permet donc pas de fusionner les échantillons M4B, Live et M4.4.

\subsection{Pourquoi le regroupement temporel n'est pas une symétrie}
\label{sec:composition_temps}
Un contre-exemple apparaît avant même d'introduire le marché. Sans bruit,
crédit ni naissances, posons $F(K)=(1-\delta)(K+AK^\gamma)$.
Deux pas successifs donnent
\[
F(F(K))=(1-\delta)^2(K+AK^\gamma)
 +(1-\delta)A\big[(1-\delta)(K+AK^\gamma)\big]^\gamma.
\]
Un pas regroupé, avec $A'=2A$ et
$1-\delta'=(1-\delta)^2$, donnerait au contraire
\[
G(K)=(1-\delta)^2(K+2AK^\gamma).
\]
En général $F\circ F\ne G$ : la seconde production se calcule sur
un stock déjà modifié. Par exemple, à $K=100$, $A=1$,
$\gamma=1/2$ et $\delta=0{,}01$, les valeurs sont environ
$118{,}142$ et $117{,}612$. Le marché ajoute une alternance de
rencontres, de contrats, d'intérêts et de faillites ; composer les lois
de Poisson ou les seuls chocs ne compose pas ces opérations couplées.

Les essais Live-v1 de la figure~\ref{fig:art_temps} mesurent la portée
collective de cette différence. Après conversion complète des paramètres,
le regroupement de deux pas donne, relativement à sa référence,
$N:1{,}067\pm0{,}018$, $K_{\rm tot}:1{,}158\pm0{,}030$ et
$Y:1{,}118\pm0{,}024$. En regroupant quatre pas, le rapport de
production vaut $1{,}346\pm0{,}026$. Le protocole part d'un état vide,
utilise trois graines et compare $1200<t_{\rm ancien}\le2000$ ;
les flux sont divisés par le facteur de regroupement, pas les stocks.
Les valeurs par graine et les autres variantes sont dans
\chemin{data_article/temps_graines.csv} et
\chemin{data_article/temps_points.csv}. Ces essais et le contre-exemple
établissent l'absence d'une symétrie exacte ; ils ne constituent pas
une étude de convergence vers une limite continue.

